\documentclass[12pt,a4paper]{article}
\usepackage[margin=1in]{geometry}
\usepackage{amsmath,amssymb}
\usepackage{hyperref}
\usepackage{booktabs}
\usepackage{graphicx}
\usepackage{natbib}

\title{\textbf{A Physically Interpretable Single-Mode Phonon-Occupation Model
for the Temperature-Dependent Zero-Field Splitting in NV Centers}}

\author{Reddouane BERRAMDANE\\
\small Independent Researcher\\
\small \href{mailto:reddoma@gmail.com}{reddoma@gmail.com}}

\date{June 2026}

\begin{document}
\maketitle

% ─────────────────────────────────────────────────────────────────────────────
\begin{abstract}
We present a compact effective model in which the temperature-dependent
part of the NV-center zero-field splitting $D(T)$ is governed by the
Bose--Einstein occupation of a single representative phonon mode.
The model is equivalent to a single-mode $\coth$ fluctuation model
after absorbing the zero-point contribution into the fitted intercept.
Calibrated to the high-precision Ouyang et al.\ dataset over 298--383\,K,
the model achieves residuals comparable to empirical polynomial fits while
retaining three physically interpretable parameters: the effective intercept
$D_0$, a coupling amplitude $A$, and an effective phonon wavenumber
$\tilde{\nu}$.
The fitted effective wavenumber,
$\tilde{\nu} = 711 \pm 48\ \mathrm{cm}^{-1}$,
lies within the diamond phonon spectrum and should be interpreted as an
effective spin--phonon energy scale rather than a resolved normal mode.
When benchmarked against the Liu et al.\ high-temperature parametrization
up to 1000\,K, the model remains within approximately 2\,MHz over the full
range ($R^2 = 0.998185$, RMSE\,=\,1.3\,MHz), while a power-law calibration
diverges outside the fitting window.
We refer to this modelling framework as the
\emph{Decoherence-Temperature Quantum Effective Model} (DTQEM),
of which the present single-mode thermal shift constitutes the
first quantitative component.
\end{abstract}

% ─────────────────────────────────────────────────────────────────────────────
\section{Introduction}

The nitrogen-vacancy (NV) center in diamond is one of the most versatile
solid-state quantum sensors \cite{Degen2017,Barry2020}.
The temperature dependence of the ground-state zero-field splitting,
$D(T)$, underpins NV-based nanothermometry \cite{Kucsko2013,Neumann2013,Toyli2013a}
and must be accounted for in precision measurements under varying thermal
conditions.
Although $D(T)$ has been measured over a wide range of temperatures
\cite{Acosta2010,Chen2011,Toyli2012,Plakhotnik2014,Ouyang2024},
the most common descriptions remain empirical polynomial fits with no
direct physical meaning.

A physical interpretation links the temperature shift to lattice expansion
and electron--phonon coupling \cite{Doherty2014}.
Recent work by Cambria et al.\ \cite{Cambria2024} showed that two effective
phonon modes spanning the acoustic and optical branches of the diamond
spectrum are sufficient to describe $D(T)$ quantitatively.
Building on this idea, we show that a single effective mode can achieve
comparable accuracy within the calibration window and strong extrapolation
performance without increasing parameter complexity.
This parsimony is attractive because it yields a model with direct physical
meaning and a controlled high-temperature asymptote.

We refer to this modelling framework as the
\textbf{Decoherence-Temperature Quantum Effective Model (DTQEM)}.
The present paper constitutes Paper~I of this series, addressing the
thermal shift of $D(T)$; Paper~II will extend the framework to
spin-lattice relaxation $1/T_1(T)$ and pure dephasing $\Gamma_\phi(T)$.

% ─────────────────────────────────────────────────────────────────────────────
\section{Model}

\subsection{Basic form}

We model the temperature-dependent zero-field splitting as
%
\begin{equation}
    D(T) = D_0 + A\,n(\tilde{\nu}, T),
    \label{eq:model}
\end{equation}
%
where
%
\begin{equation}
    n(\tilde{\nu}, T) =
    \frac{1}{\exp\!\left(\dfrac{hc\tilde{\nu}}{k_B T}\right) - 1}
    =
    \frac{1}{\exp\!\left(\dfrac{\tilde{\nu}}{k_{B,\mathrm{cm}}\,T}\right) - 1}
    \label{eq:BE}
\end{equation}
%
is the Bose--Einstein occupation number, $\tilde{\nu}$ is the effective
phonon wavenumber (cm$^{-1}$), and
$k_{B,\mathrm{cm}} = k_B/hc = 0.69503476\ \mathrm{cm}^{-1}\,\mathrm{K}^{-1}$
is the Boltzmann constant expressed in wavenumber units.
Here $D_0$ is an effective low-temperature intercept, $A$ is a coupling
amplitude, and $\tilde{\nu}$ is an effective phonon energy scale.

\subsection{Relation to the \texorpdfstring{$\coth$}{coth} fluctuation form}

Although the Bose--Einstein occupation number $n(\tilde{\nu},T)$
appears naturally as a thermal shift term, it is instructive to note
its connection to the more familiar phonon-fluctuation form.
The mean-square displacement of a harmonic oscillator mode scales as
%
\begin{equation}
    \langle Q^2 \rangle \propto 2n(\tilde{\nu},T)+1
    = \coth\!\left(\frac{hc\tilde{\nu}}{2k_BT}\right),
\end{equation}
%
so a general linear coupling gives
%
\begin{equation}
    D(T) = D_{\mathrm{bare}} + \lambda\,[2n(\tilde{\nu},T)+1].
\end{equation}
%
Since the zero-point term $\lambda$ is temperature-independent,
it can be absorbed into the intercept, yielding
%
\begin{equation}
    D(T) = \underbrace{(D_{\mathrm{bare}}+\lambda)}_{D_0}
           + \underbrace{2\lambda}_{A}\,n(\tilde{\nu},T).
\end{equation}
%
The present model~\eqref{eq:model}--\eqref{eq:BE} is therefore
\emph{equivalent to a single-mode $\coth$ fluctuation model},
with the zero-point contribution absorbed into the fitted intercept $D_0$.
Consequently, $D_0$ should be interpreted as an effective low-temperature
intercept that includes both the bare zero-field splitting and the
zero-point lattice contribution, rather than as a precise cryogenic
value of $D(0)$.

At high temperature ($k_BT \gg hc\tilde{\nu}$), $n \approx k_BT/hc\tilde{\nu}$
and $D(T)$ becomes approximately linear in $T$.
At low temperature, $D(T)$ approaches $D_0$ exponentially.
The single-mode ansatz is adopted here on grounds of parsimony: it achieves
strong fit quality within the calibration window while preserving a
physically interpretable structure.

% ─────────────────────────────────────────────────────────────────────────────
\section{Results}

\subsection{Calibration on high-precision data}

We fit the model to the dataset of Ouyang et al.\ \cite{Ouyang2024},
taken from Table~1 of that work
(18 points, 298.15--383.15\,K, uncertainties
$u(D) = (3.6\text{--}8.5)\times10^{-5}$\,GHz).
The fitted parameters are
%
\begin{align}
    D_0    &= 2.8764\ \mathrm{GHz}, \\
    A      &= -0.1889 \pm 0.0299\ \mathrm{GHz}, \\
    \tilde{\nu} &= 711 \pm 48\ \mathrm{cm}^{-1}.
\end{align}
%
The corresponding fit quality is $R^2 = 0.999946$ with
RMSE\,$= 1.75\times10^{-5}$\,GHz.

The Gaussian approximation to the parameter posterior is justified by
the small relative uncertainties:
$\sigma(D_0)/D_0 < 0.02\%$,
$\sigma(A)/|A| \approx 16\%$, and
$\sigma(\tilde{\nu})/\tilde{\nu} \approx 7\%$.

The extracted $D_0$ is slightly lower than the commonly cited
$D(0) \approx 2.8775$--$2.8780$\,GHz \cite{Doherty2014}.
This difference is expected: $D_0$ is an effective intercept that
absorbs the zero-point lattice contribution (see Section~2), and should
not be compared directly to cryogenic measurements of $D(0)$ without
accounting for that offset.

\subsubsection*{Parameter estimation and uncertainty propagation}

Model parameters were estimated by weighted nonlinear least squares
using \texttt{scipy.optimize.curve\_fit} with the reported standard
uncertainties $u(D)$ of Ouyang et al.\ as $1\sigma$ weights
(\texttt{absolute\_sigma=True}), yielding the parameter covariance
matrix $\Sigma$ directly in physical units.
The $1\sigma$ parameter uncertainties are obtained as
$\sigma_i = \sqrt{\Sigma_{ii}}$.

To propagate parameter uncertainty to the model predictions, we drew
$N = 10\,000$ independent samples from the joint Gaussian posterior
$\mathcal{N}(\hat{p}, \Sigma)$, evaluated $D(T)$ for each sample over
the full extrapolation range 280--1020\,K, and computed pointwise
2.5th and 97.5th percentiles.
The resulting shaded regions in Fig.~2(a) represent
\emph{95\% parameter-uncertainty bands}; they reflect only the
propagated uncertainty in the fitted parameters and do not include
measurement noise at the extrapolation temperatures.
The same procedure was applied independently to the power-law model
to enable a direct visual comparison.
The Liu et al.\ values used for residual evaluation were digitized
from Fig.~2(c) of Ref.~\cite{Liu2019} and carry no independently
assigned measurement uncertainty; the $\pm5$\,MHz shading in
Fig.~2(b) is a visual guide only.

\subsection{Model comparison}

All three models achieve comparable fit quality within the calibration
window (Table~\ref{tab:comparison}).
The main distinction is interpretability: the present model uses
physically meaningful parameters, whereas the polynomial coefficients
are purely empirical.

\begin{table}[h]
\centering
\caption{Model comparison in the 298--383\,K calibration window.}
\label{tab:comparison}
\begin{tabular}{lccc}
\toprule
Model & $R^2$ & RMSE (GHz) & Parameters \\
\midrule
This work (DTQEM)  & 0.999946 & $1.75\times10^{-5}$ & 3 \\
Ouyang polynomial  & 0.999932 & $1.95\times10^{-5}$ & 3 \\
Power law          & 0.999930 & $1.98\times10^{-5}$ & 3 \\
\bottomrule
\end{tabular}
\end{table}

\begin{figure}[h]
    \centering
    \includegraphics[width=0.95\textwidth]{paper/fig1_dtqem_fit_final.pdf}
    \caption{(a) Fit to the Ouyang et al.\ data. (b) Fit residuals. (c) $R^2$ comparison. All three models achieve comparable fit quality; RMSE values are below the reported measurement uncertainty.}
    \label{fig:fit}
\end{figure}

\subsection{Extrapolation benchmark: comparison with Liu et al.\ up to 1000\,K}

We apply the parameters obtained from the Ouyang calibration directly
to the high-temperature $D(T)$ parametrization of Liu et al.\ \cite{Liu2019}.
The values were sampled from the published curve using digital extraction
and are used here as a benchmark.

The deviation remains below 2\,MHz across the full range (Table~\ref{tab:liu}).
In contrast, a power-law model exhibits deviations that grow rapidly
outside the fitting window.
Over the full 300--1000\,K range, the DTQEM extrapolation yields
$R^2 = 0.998185$ and an RMSE of 1.32\,MHz.
The Liu et al.\ parametrization lies within the DTQEM 95\%
parameter-uncertainty band across the full 300--1000\,K range,
whereas the power-law band diverges above approximately 600\,K.

\begin{figure}[h]
    \centering
    \includegraphics[width=0.95\textwidth]{paper/fig2_MC_bands_final.pdf}
    \caption{Extrapolation of the DTQEM and power-law models from the 298--383\,K calibration window (gray shading) to the Liu et al.\ parametrization up to 1000\,K. (a) Model predictions with 95\% parameter-uncertainty bands obtained by Monte Carlo sampling ($N=10\,000$). (b) Residuals with respect to the digitized Liu et al.\ values; the $\pm5$\,MHz gray band is a visual guide.}
    \label{fig:extrapolation}
\end{figure}

\begin{table}[h]
\centering
\caption{Comparison with the Liu et al.\ parametrization up to 1000\,K.}
\label{tab:liu}
\begin{tabular}{cccc}
\toprule
$T$ (K) & $D_{\mathrm{Liu}}$ (GHz) & $D_{\mathrm{model}}$ (GHz) & Difference (MHz) \\
\midrule
300  & 2.8680 & 2.8699 & $-1.9$ \\
400  & 2.8600 & 2.8605 & $-0.5$ \\
500  & 2.8490 & 2.8484 & $+0.6$ \\
600  & 2.8360 & 2.8344 & $+1.6$ \\
700  & 2.8210 & 2.8194 & $+1.6$ \\
800  & 2.8050 & 2.8035 & $+1.5$ \\
900  & 2.7880 & 2.7872 & $+0.8$ \\
1000 & 2.7700 & 2.7704 & $-0.4$ \\
\bottomrule
\end{tabular}
\end{table}

% ─────────────────────────────────────────────────────────────────────────────
\section{Discussion}

\subsection{Physical interpretation}

The fitted wavenumber $\tilde{\nu} = 711 \pm 48\ \mathrm{cm}^{-1}$
($\approx 88\ \mathrm{meV}$) lies between the acoustic-phonon regime and
the zone-center optical phonon at $1332\ \mathrm{cm}^{-1}$, and is
bracketed by the two effective modes reported by Cambria et al.\
\cite{Cambria2024}.
It should therefore be interpreted as an effective energy scale
representing the dominant spin--phonon coupling channel,
not as a spectroscopically resolved single lattice mode.

A Debye-integral model was also tested as an alternative,
replacing $n(\tilde{\nu}, T)$ with the integrated phonon density of
states weighted by a Debye cutoff frequency $\tilde{\nu}_D$.
Although physically motivated, this model requires numerical integration,
introduces an additional free parameter, and yields no measurable
improvement in $R^2$ or RMSE within the calibration window.
The single-mode ansatz is therefore retained on grounds of parsimony.

\subsection{Identifiability and scope}

The model should be interpreted as a compact effective description whose
individual parameters are constrained by the present dataset but not
fully uniquely identifiable without broader temperature coverage.
The agreement with the Liu et al.\ parametrization indicates a robust
thermal trend rather than a fitting artifact.
The model should be used with caution above 1000\,K, where higher-order
phonon--phonon interactions and anharmonic effects become significant.

The uncertainty on the effective phonon wavenumber,
$\sigma(\tilde{\nu}) = 48\ \mathrm{cm}^{-1}$ (${\sim}7\%$),
could be substantially reduced by incorporating low-temperature
measurements (e.g., $T < 100$\,K), where the exponential sensitivity
of $n(\tilde{\nu}, T)$ to $\tilde{\nu}$ is strongest.
Alternatively, a joint fit to both $D(T)$ and $1/T_1(T)$ over a common
temperature range would provide independent constraints on the coupling
amplitude $A$ and the effective wavenumber $\tilde{\nu}$, improving
parameter identifiability without requiring new instrumentation beyond
what is standard in NV-center spectroscopy.

\subsection{Implications for thermometry}

The model combines a physical low-temperature limit, a controlled
high-temperature asymptote, and strong extrapolation performance,
making it preferable to unconstrained polynomial fits for NV
thermometry applications.

\subsection{Possible extension to relaxation and dephasing (Paper II)}

The same phonon-occupation framework may be extended to spin-lattice
relaxation rates $1/T_1(T)$ and pure dephasing $\Gamma_\phi(T)$.
For pure dephasing, the natural ansatz is:
%
\begin{equation}
    \Gamma_\phi(T) = \Gamma_0 + \alpha\,[2n(\tilde{\nu}_\phi, T)+1]
    = \Gamma_0 + \alpha\,\coth\!\left(\frac{hc\tilde{\nu}_\phi}{2k_BT}\right),
\end{equation}
%
while for spin-lattice relaxation:
%
\begin{equation}
    \frac{1}{T_1(T)} \propto n(\tilde{\nu},T)\,[n(\tilde{\nu},T)+1].
\end{equation}
%
A key question for Paper~II is whether the same effective wavenumber
$\tilde{\nu}$ governs both $D(T)$ and $\Gamma_\phi(T)$
(single dominant spin--phonon channel), or whether distinct modes must
be introduced.
A robust quantitative test would require experimental data spanning
at least 100--600\,K with sub-percent precision on $T_1$.
This is left as a concrete objective for Paper~II of this series.

% ─────────────────────────────────────────────────────────────────────────────
\section{Conclusion}

We have presented a physically motivated single-mode occupation-number
model for $D(T)$ in NV centers, and shown its equivalence to a
single-mode $\coth$ fluctuation model with an absorbed zero-point term.
The model matches the best published polynomial fit within the reported
experimental uncertainties of the Ouyang et al.\ dataset while providing
direct physical interpretability.
The extracted effective wavenumber
$\tilde{\nu} = 711 \pm 48\ \mathrm{cm}^{-1}$ is consistent with the
diamond phonon spectrum.
Most importantly, when calibrated only on 298--383\,K, the model
reproduces the Liu et al.\ parametrization up to 1000\,K with
$R^2 = 0.998185$, whereas a power-law calibration performs poorly
outside the fitting window.
These results establish DTQEM as a compact, physically interpretable
alternative to purely empirical descriptions of $D(T)$ and lay the
groundwork for a unified spin--phonon framework connecting thermal
shifts and decoherence rates in NV centers.

% ─────────────────────────────────────────────────────────────────────────────
\section*{Data Availability}

The Ouyang et al.\ data used for model calibration are available in
Table~1 of Ref.~\cite{Ouyang2024}.
The Liu et al.\ parametrization values used for the extrapolation
benchmark were digitally extracted from Fig.~2(c) of Ref.~\cite{Liu2019}
and are tabulated in the present manuscript.
A preprint version of this work, along with the Python fitting code
and the digitized benchmark values, is available at
DOI: \texttt{[Zenodo DOI upon acceptance]}.

% ─────────────────────────────────────────────────────────────────────────────
\section*{Acknowledgments}

The author thanks colleagues and collaborators for valuable discussions
and methodological feedback that improved the clarity and rigor of the
manuscript.
The author also acknowledges the assistance of DeepSeek, Claude, Grok,
and Perplexity AI systems for help with code development, data extraction,
literature survey, \LaTeX\ manuscript preparation, and scientific
discussion.

% ─────────────────────────────────────────────────────────────────────────────
\section*{Conflict of Interest}

The author declares no conflict of interest.

% ─────────────────────────────────────────────────────────────────────────────
\bibliographystyle{unsrt}
\begin{thebibliography}{99}

\bibitem{Degen2017}
C.~L. Degen, F.~Reinhard, and P.~Cappellaro,
\textit{Rev.\ Mod.\ Phys.}\ \textbf{89}, 035002 (2017).

\bibitem{Barry2020}
J.~F. Barry et al.,
\textit{Rev.\ Mod.\ Phys.}\ \textbf{92}, 015004 (2020).

\bibitem{Kucsko2013}
G.~Kucsko et al.,
\textit{Nature}\ \textbf{500}, 54 (2013).

\bibitem{Neumann2013}
P.~Neumann et al.,
\textit{Nano Lett.}\ \textbf{13}, 2738 (2013).

\bibitem{Toyli2013a}
D.~M. Toyli et al.,
\textit{Proc.\ Natl.\ Acad.\ Sci.\ USA}\ \textbf{110}, 8417 (2013).

\bibitem{Acosta2010}
V.~M. Acosta et al.,
\textit{Phys.\ Rev.\ Lett.}\ \textbf{104}, 070801 (2010).

\bibitem{Chen2011}
X.-D. Chen et al.,
\textit{Appl.\ Phys.\ Lett.}\ \textbf{99}, 161903 (2011).

\bibitem{Toyli2012}
D.~M. Toyli et al.,
\textit{Phys.\ Rev.\ X}\ \textbf{2}, 031001 (2012).

\bibitem{Plakhotnik2014}
T.~Plakhotnik et al.,
\textit{Nano Lett.}\ \textbf{14}, 4989 (2014).

\bibitem{Ouyang2024}
K.-C. Ouyang et al.,
\textit{Phys.\ Rev.\ Research}\ \textbf{6}, 023078 (2024).

\bibitem{Doherty2014}
M.~W. Doherty et al.,
\textit{Phys.\ Rev.\ B}\ \textbf{90}, 041201(R) (2014).

\bibitem{Cambria2024}
M.~C. Cambria et al.,
arXiv:2404.08388v2 (2024).

\bibitem{Jarmola2012}
A.~Jarmola et al.,
\textit{Phys.\ Rev.\ Lett.}\ \textbf{108}, 197601 (2012).

\bibitem{Liu2019}
G.-Q. Liu et al.,
\textit{Nat.\ Commun.}\ \textbf{10}, 1344 (2019).

\end{thebibliography}

\end{document}
